\section{Mode F1 -- Production normale (couplé à A2)}
\label{secModeF1}

\subsection{Conditions d'entrée}

Le mode F1 est atteint depuis le mode A1, sur demande de marche automatique
formulée par l'opérateur.

\subsection{Comportement attendu}

L'objectif de ce mode est le transfert, par le robot, du gobelet porté par
une palette pleine arrivée dans le guichet, vers la prochaine palette vide
détectée dans le guichet.

\begin{itemize}
    \item \textbf{Convoyeur / guichet} : la \textbf{capacité du guichet
    doit rester unitaire} en toute circonstance : une palette ne peut
    entrer dans le guichet que si celui-ci est libre. Le maintien en poste
    d'une palette doit permettre au robot d'y effectuer la saisie ou la
    dépose d'un gobelet lorsque la palette est immobilisée.
    \item \textbf{Robot} : le robot doit réaliser la saisie du gobelet sur
    la palette pleine présente dans le guichet, puis sa dépose sur la
    prochaine palette vide détectée dans le guichet. Le transfert engagé
    doit toujours être mené à son terme, sans interruption prématurée par
    un changement de mode.
    \item \textbf{IHM} : l'état du guichet doit être signalé à l'opérateur
    à chaque étape du cycle (poste libre, palette en cours d'arrivée,
    palette en poste avec gobelet à traiter, palette en cours de
    libération), par une signalisation dédiée et distincte à chaque étape.
\end{itemize}

\begin{UPSTIwarning}{Contrainte de capacité unitaire du guichet}
Comme rappelé dans le TD de conception, le guichet fonctionne comme un
\textbf{sas} : sa capacité est unitaire. Cette contrainte, déjà valable en
fonctionnement du guichet seul, reste impérative en mode F1, y compris
pendant les phases où le robot manipule un gobelet.
\end{UPSTIwarning}

\subsection{Articulation avec le mode A2}

Le mode F1 est couplé à un mode d'arrêt A2 (\emph{arrêt en fin de cycle}) :
une demande d'arrêt formulée par l'opérateur pendant que le robot est en
train de transférer un gobelet ne doit pas interrompre ce transfert en
cours ; l'installation ne rejoint le mode A2 qu'une fois le transfert de
gobelet en cours effectivement terminé.

\subsection{Conditions de sortie}

\begin{itemize}
    \item Vers le mode \textbf{A2}, sur demande d'arrêt de l'opérateur, une
    fois le transfert de gobelet en cours terminé.
    \item Retour vers le mode \textbf{F1} depuis A2, dès que la demande
    d'arrêt n'est plus active.
    \item Depuis A2 vers le mode \textbf{A1}, en fin de cycle, une fois la partie opérative en conditions initiales.q
\end{itemize}
